\chapter{Пропагация убеждений как лавина}
\label{ch:propagation}

\section{Постановка: затухающий градиент уверенности}

Когда в память приходит новый факт $A$, он линкуется к начальным узлам (например, к
последним $k$ активным), образуя окрестность $\Neigh(A)$. Мы линкуем $A\to\Neigh(A)$
и смотрим, как меняется уверенность соседей. Если изменение велико, распространяем
дальше, $\Neigh(A)\to\Neigh(\Neigh(A))$, — как затухающий градиент. Знание, ни на
что не повлиявшее, чуть пошатнёт соседей и затухнет; знание с большими градиентами
вызовет \emph{лавину}. Большую лавину называем «кризисом мировоззрения» — и такое
бывает с настоящим умом.

\section{Это самоорганизованная критичность}

У картины есть точное имя. «Лавина» — это \emph{нейронная лавина}, а система у
порога между затуханием и разрастанием — \emph{самоорганизованная критичность}
(Bak–Tang–Wiesenfeld; для мозга — Beggs–Plenz): каскады следуют степенному закону.

Ключевой управляющий параметр — не «ручная константа затухания», а \emph{ветвящийся
коэффициент} $\sigma$ (среднее число узлов, активируемых одним активным):
\[
\sigma<1:\ \text{субкритично};\quad \sigma=1:\ \text{критично};\quad
\sigma>1:\ \text{суперкритично (runaway)}.
\]
Для ветвящегося процесса ожидаемый размер лавины конечен при $\sigma<1$ и равен
\begin{equation}
\mathbb{E}[\text{размер}]\;=\;\frac{1}{1-\sigma}.
\label{eq:avalanche}
\end{equation}
Держа $\sigma$ чуть меньше единицы, система почти всегда сама гасит распространение,
но изредка пропускает крупную переоценку. Почему это лучше фиксированной константы:
константа не адаптируется к плотности графа (в плотном регионе та же константа даёт
лавину, в разреженном — глушит всё). Нормировка под целевое $\sigma$ адаптивна.

\begin{intuition}[Куча песка]
Классический образ SOC — куча песка, на которую сыплют по песчинке. Почти каждая
не делает ничего; изредка одна вызывает оползень; ещё реже — обвал, перестраивающий
склон. Куча сама держится у критического угла. Убеждения — так же: почти любой факт
проходит без следа, изредка меняет соседние взгляды, совсем изредка обрушивает
пласт мировоззрения. Формула \eqref{eq:avalanche} говорит: чем ближе $\sigma$ к
единице, тем крупнее (и реже) обвалы.
\end{intuition}

\begin{pitfall}[SOC — диагностическая цель, а не доказанный закон]
Самоорганизованная критичность соблазнительна и красиво объясняет редкие «кризисы
мировоззрения». Но на \emph{типизированном, знаковом, направленном и меняющемся}
графе памяти она — \emph{рамка}, а не выведенный закон: строгих гарантий, что
пропагация убеждений действительно самонастраивается к критичности, у нас нет.
Поэтому используем её \textbf{инженерно, как цель мониторинга, а не как механизм,
который «строим»}: меряем распределение размеров ревизий (лавин) и держим его в
здоровом диапазоне. Всё глохнет — субкритично (память ригидна, не обновляется); всё
взрывается — суперкритично (нестабильность, «зависает на сутки»); изредка бывают
крупные ревизии, обычно мелкие — здорово. Ветвящийся коэффициент $\sigma$ — это
управляющая ручка, которой мы держим систему субкритичной, а не утверждение о том,
что «так оно и есть». Конкретное правило-предохранитель: мерим $\sigma$ по телеметрии
и при $\sigma\gtrsim 1{,}2$ \textbf{программно врубаем жёсткое затухание}. Иначе
получишь суперкритический разгон — бесконечный цикл обновления уверенностей, который
съест LLM-слоты и всё подвесит (гл.~\ref{ch:runtime}). Не доказываем критичность
аналитически — держим субкритичность рычагом.
\end{pitfall}

\begin{openq}[Порог $\sigma$ — заглушка первой версии]
Значение $\sigma\approx 1{,}2$ — \emph{инженерная заглушка для первой реализации},
калибруется на стенде. \emph{Почему такой порядок, а не 1,01 или 2}: чуть выше
единицы, чтобы крупные ревизии случались, но \emph{редко} (само явление «кризиса
мировоззрения» должно быть редким); слишком близко к 1 (1,01) — крупные каскады
зачастят; заметно выше (2) — суперкритический разгон. То есть «чуть-чуть за порогом,
с запасом» — отсюда $1{,}1$--$1{,}3$; конкретика — со стенда. И само $\sigma$ на циклическом типизированном графе ещё надо
\emph{операционализировать}: что именно считать «одним активированным потомком», в
каком окне, как учитывать повторные визиты, разные типы рёбер и уже активные узлы.
Идея (держать субкритичность) верна; метрика требует доводки.
\end{openq}

\section{Где вся сложность: локальное правило}

Пороговая логика — лёгкая часть (процентов двадцать). Восемьдесят процентов
спрятаны в локальном правиле $\Delta\conf(B\mid A)$:
\begin{equation}
\Delta\conf(B\mid A)\;=\;\eta\;\cdot\;\underbrace{\mathrm{sign}(r_{AB})}_{(1)}\;\cdot\;
\underbrace{w_{AB}}_{}\;\cdot\;\underbrace{c_A\,\mathrm{trust}(A)}_{(2)}\;\cdot\;
\underbrace{\iota(A,B)}_{(3)} .
\label{eq:deltaconf}
\end{equation}
Три множителя, без которых правило сломано.

\paragraph{(1) Знак.} Новый факт сосед либо \emph{подтверждает}, либо
\emph{опровергает}. Беззнаковое распространение не отличит «всё усилилось»
(озарение) от «всё перевернулось» (кризис). Нужны знаковые/типизированные рёбра.

\paragraph{(2) Гейт по надёжности $A$ (карантин).} Новый факт может сам быть ложным.
Множитель $c_A\,\mathrm{trust}(A)$ вводит \emph{карантин}: новый факт пропагирует
слабо (малый эффективный $\sigma$), пока сам не заработал уверенность через
корроборацию или выживание в скрабе. Иначе один ненадёжный вход вызовет кризис
мировоззрения. По сути — байесовское обновление с правдоподобием, взвешенным
доверием к источнику.

\paragraph{(3) Дисконт независимости $\iota(A,B)$.} Граф циклический, и по петле $A$
может подтвердить сам себя (\emph{эхо-камера}, петлевое распространение), раздувая
уверенность из ничего. Дисконтируем по общей родословной свидетельств (провенанс), а
не считаем активации: число упоминаний $\ne$ число независимых свидетельств. Это
ровно механизм человеческих эхо-камер.

\begin{pitfall}[Дисконт независимости — скрытая $O(|V|{+}|E|)$-ловушка; чините сигнатурами]
Наивно проверять независимость $A$ и $B$ — значит сравнивать их \emph{деревья
провенанса} (наименьший общий предок / обход), а это в худшем случае
$O(|V|{+}|E|)$. Запускать такое \emph{на каждый} обновляемый узел во время лавины —
и «миллисекунды» превращаются в секунды. Оптимизация: для узлов держать
\emph{сигнатуру родословной провенанса} — фильтр Блума или MinHash. Тогда проверка
сводится к побитовому AND / перекрытию сигнатур за $O(1)$:
\[
\iota(A,B)\;\approx\;1-\mathrm{overlap}\big(\mathrm{sig}(A),\ \mathrm{sig}(B)\big).
\]
Честный нюанс: Bloom даёт ложные срабатывания $\Rightarrow$ иногда \emph{пере-}
дисконтируем (сочтём независимые источники общими). Но это \emph{безопасная} сторона
ошибки — лучше скромнее считать свидетельства, чем раздуть уверенность. Эта же
сигнатура потом переиспользуется при разрешении противоречий во сне
(глава~\ref{ch:sleep}).

\emph{Когда что из двух}: Bloom отвечает на \emph{булев} вопрос «есть ли у $A$ и $B$
\emph{хоть один} общий предок-источник» — этого достаточно для грубого гейта
«независимы / нет». MinHash оценивает \emph{долю} перекрытия родословных (Jaccard) и
нужен, когда важна \emph{степень} зависимости — то есть плавный дисконт $\iota\in[0,1]$,
а не «да/нет». Правило: начинаем с Bloom (дешевле, меньше памяти); переходим к MinHash
там, где грубого гейта мало и хочется градуированного веса свидетельств.
\end{pitfall}

\section{Пороговая логика — это гистерезис}

Два порога $\theta_{\text{emergent}}>\theta_{\text{temp}}>0$:
\begin{itemize}
\item $\Delta\conf>\theta_{\text{emergent}}$ — обязательное распространение «в
ширину» (крупная переоценка);
\item $\theta_{\text{emergent}}>\Delta\conf>\theta_{\text{temp}}$ — узлы фронта
помечаются \texttt{need\_to\_revise};
\item $\Delta\conf<\theta_{\text{temp}}$ — распространение глохнет здесь.
\end{itemize}
Два порога вместо одного — \emph{гистерезис} (триггер Шмитта): нет «дребезга» у
единственного порога. Тонкость: помеченные во второй зоне узлы позже
переинспектируются ночным обходчиком, и старые записи, чью уверенность мы не
трогали, могут «потечь обратно» — такое бывает и в жизни.

\section{Пороги растут с глубиной консолидации}

Пороги — не глобальные константы. $\theta_{\text{emergent}}(v)$ растёт с глубиной
консолидации узла $v$: периферийный узел перевернуть дёшево, а identity-core в слое
«жизнь» — дорого, нужно больше свидетельств, чтобы вообще запустить лавину. Это и
есть ось \emph{пластичность\,--\,упрямство}: чем глубже убеждение продиффундировало
в медленные слои, тем выше эффективный порог опрокидывания. Одна ручка — и она же
задаёт, насколько Рин склонна к belief perseverance. Псевдокод пропагации с
контролем $\sigma$ — в приложении~\ref{ch:algorithms}.
